%%%%%%%%%%%%%%%%%%%%%%%%%%%%%%%%%%%%%%%%%%%%%%%%%%%%%%%%%%%
\chapter*{Conclusion générale}
\addcontentsline{toc}{chapter}{Conclusion générale}
%%%%%%%%%%%%%%%%%%%%%%%%%%%%%%%%%%%%%%%%%%%%%%%%%%%%%%%%%%%

Le projet \textbf{Certif-fun} a abouti à la création d'un écosystème innovant alliant architecture web distribuée et intelligence artificielle pour répondre au défi de l'intégrité des certifications à distance.

\paragraph{Récapitulation des objectifs :} 
Initialement centré sur la sécurisation des évaluations, le projet a relevé le défi de l'automatisation de la surveillance. Les objectifs techniques (microservices, scalabilité) et fonctionnels (proctoring temps réel, génération de quiz) ont été validés, offrant une solution complète et opérationnelle.

\paragraph{Contributions et résultats :} 
Les réalisations majeures se résument à :
\begin{itemize}[noitemsep, topsep=1pt]
    \item \textbf{Surveillance intelligente} : Intégration de \textbf{YOLOv8} et \textbf{MediaPipe} pour un suivi autonome du regard et des objets suspects.
    \item \textbf{Pédagogie augmentée} : Génération automatique de quiz via \textbf{Ollama}, garantissant la confidentialité des données.
    \item \textbf{Agilité DevOps} : Mise en œuvre d'un pipeline \textbf{CI/CD} performant assurant des déploiements fluides.
\end{itemize}

\paragraph{Limites du travail :} 
Le système rencontre toutefois des contraintes matérielles liées au coût computationnel de l'IA sur un serveur unique. La précision reste également sensible à l'environnement de l'apprenant (éclairage, caméra), et l'authentification mériterait d'être renforcée par un contrôle biométrique continu.

\paragraph{Perspectives futures :} 
Les pistes d'évolution incluent la migration vers \textbf{Kubernetes} pour une scalabilité élastique, l'intégration de la \textbf{reconnaissance faciale continue} et le déport de l'IA vers le navigateur (\textit{Edge Computing}). L'ajout d'analyses prédictives permettrait également d'anticiper les risques de décrochage.

En conclusion, \textbf{Certif-fun} constitue un socle technologique solide, démontrant que l'IA peut devenir le garant de la crédibilité des diplômes numériques de demain.
